\documentclass[preprint,12pt,authoryear]{elsarticle}

\usepackage{iftex}
\ifTUTeX                          % XeTeX/LuaTeX (tectonic): Unicode fonts, so ı İ ş ğ can be typed directly
  \usepackage{fontspec}           % Latin Modern by default
\else                             % pdfLaTeX
  \usepackage[T1]{fontenc}
  \usepackage[utf8]{inputenc}
  \usepackage{lmodern}
\fi
\usepackage{amsmath,amssymb}
\usepackage{booktabs}
\usepackage{array}
\usepackage{microtype}
\usepackage{graphicx}
\usepackage[hidelinks]{hyperref}

\graphicspath{{figures/}}

% A figure or table that exists only once its results have been run.
\newcommand{\maybefigure}[2]{\IfFileExists{figures/#1.pdf}{\includegraphics[width=#2]{#1}}{\fbox{\parbox{0.9\linewidth}{\centering\small figures/#1.pdf not made yet (\texttt{make figures})}}}}
\newcommand{\maybetable}[1]{\IfFileExists{figures/#1.tex}{\input{figures/#1.tex}}{\fbox{\small figures/#1.tex not made yet (\texttt{make figures})}}}

\journal{(journal to be decided)}

\begin{document}

\begin{frontmatter}

\title{Generating virtual GNSS stations from their neighbours with a conditional diffusion model: design, baselines and a first model\tnoteref{wr}}
\tnotetext[wr]{Working report, drafted with Claude Code. It records the state of the project and is updated as results come in.}

\author[ytu]{H. Oğuz Bolat\corref{cor}}
\ead{oguz.bolat@std.yildiz.edu.tr}
\cortext[cor]{Corresponding author}
\affiliation[ytu]{organization={Department of Geomatics Engineering, Yıldız Technical University},
  city={Istanbul}, country={T\"urkiye}}

\begin{abstract}
Continuous GNSS networks are dense in a few regions and sparse in most others, yet many applications need displacement time series at locations without a station. We propose to generate daily residual series at arbitrary locations from the stations around them with a conditional diffusion model, and to judge the generated series as samples rather than as point predictions. Because a station's series is only partly predictable from its neighbours, a regression model returns series that are too smooth; a generative model can reproduce the station-specific noise while inheriting the shared signal. The model is trained and validated on 1,251 stations in and around California, where real stations can be hidden and their generated counterparts compared with the truth, and tested on the network of T\"urkiye, whose median nearest-neighbour distance (60~km) is about six times California's. To learn sparse geometries from a dense network, training samples are drawn from networks thinned by minimum-separation thinning, which, unlike random subsampling, reproduces the station spacing of T\"urkiye. We describe the data, the sampling design and the evaluation protocol, which combines point, probabilistic and realism scores with planted fault-slip events, and report three baselines against which the diffusion model will be measured. On held-out stations, the short-period error of a weighted neighbour mean depends little on network density, because the shared short-period signal is dominated by common-mode error coherent over hundreds of kilometres. Adding white and flicker noise fitted to the neighbours yields a well-calibrated generative baseline, with a continuous ranked probability score 26\% below that of the mean alone, but its samples do not reproduce the decrease of inter-station correlation with distance. The diffusion model lowers the score by a further 3--4\% at every density and reproduces the decrease of correlation with distance to within 0.015 in most distance classes, Kriging and a space--time Gaussian process fitted to the training years reproduce the decrease of correlation with distance as well, and the best of them comes within 3\% of the diffusion model's score. Pooled over stations the diffusion model's short-period intervals are calibrated; but, like the baseline, it cannot tell from the neighbours how noisy the station at a location is, so its ensembles are too wide for quiet stations and too narrow for noisy ones, which makes them appear too wide at the typical station. Given the station's noise level as an optional input, a second model is calibrated across stations of every noise level and loses most of its excess variability. Planted fault-slip events reach a virtual station only when they are larger than the gaps between stations. Trained on ordinary years, the diffusion model passes them on no better than a weighted mean. Trained with simple synthetic transients planted in its inputs, it recovers a median of 52\% of the planted signal on the full network, close to the 56\% of a space--time Gaussian process, without losing accuracy on ordinary days; and at the real 2019 Ridgecrest earthquakes it reproduces the offsets of stations 55--175~km from the rupture from a network thinned to 60~km with a median error of 3.5~mm, where every method trained without transients errs by 9~mm or more.
\end{abstract}

\begin{keyword}
GNSS time series \sep virtual stations \sep network densification \sep diffusion models \sep probabilistic evaluation
\end{keyword}

\end{frontmatter}

\section{Introduction}
\label{sec:intro}

Daily position time series from continuous GNSS stations record interseismic loading, coseismic offsets, postseismic relaxation, slow slip and environmental loading. The Nevada Geodetic Laboratory (NGL) processes more than 17,000 stations worldwide with a uniform strategy \citep{blewitt2018,ngl2026}, but their distribution is very uneven: in California the median distance between a station and its nearest neighbour is 9~km, whereas in T\"urkiye, one of the most seismically active countries in the world, it is 60~km. Deformation with a footprint smaller than the station spacing, such as shallow slow slip or afterslip on a single fault segment, is then recorded by a few stations at most, or not at all.

This paper asks whether the series that a station \emph{would} record at a location without one can be generated from the stations around it. Spatial interpolation of displacement is classical, and graph neural networks have been used to estimate time series at unobserved locations of sensor networks \citep{wu2021ignnk} and to impute missing observations \citep{cini2022grin,marisca2022spin}. These methods return a single prediction. For GNSS series this is a limitation in principle rather than in degree: a station's daily residual is the sum of a part shared with its neighbours, namely common-mode error from orbits, clocks and the reference frame, environmental loading and tectonic transients \citep{wdowinski1997,dong2006,kreemer2021,li2025loading}, and a part that is its own, namely monument instability, multipath and other site effects, whose temporally correlated noise is well documented \citep{williams2004,langbein2008}. The second part cannot be predicted from other stations. Predicting a station from its neighbours is also the basis of recent network filtering, where the prediction is subtracted as noise \citep{bachelot2025}; there, the unpredictable remainder is the desired output. Here it is the part to be reproduced.

In the benchmark that preceded this work, the best deterministic predictor from neighbours reduced the median normalised error of a held-out station's short-period variations only from 0.73, for predicting zero, to about 0.52. A regression model would therefore generate virtual stations that are systematically too smooth and too quiet, with too little variance and the wrong spectrum. A conditional generative model can instead draw the station-specific part from a learned distribution while taking the shared part from the neighbours. We use a conditional diffusion model \citep{ho2020ddpm,karras2022edm}, whose samples can be scored against the truth with proper scoring rules \citep{gneiting2007} and whose realism can be checked against the spectra and spatial correlations of real stations.

The study is designed around network densification. The model is trained and validated in California, where 173 stations are held out of training entirely and their generated counterparts compared with their real series, and it is tested without any tuning in T\"urkiye. Section~\ref{sec:related} reviews related work, Section~\ref{sec:data} describes the data, Section~\ref{sec:method} the problem, the sampling design, the model and the baselines, and Section~\ref{sec:eval} the evaluation protocol. Section~\ref{sec:results} reports the results, and Section~\ref{sec:status} the decisions taken so far and the work that remains.

\section{Related work}
\label{sec:related}

\subsection{Shared and station-specific signal in GNSS series}

A station's daily residuals combine signal shared with its region and signal of its own. Spatial filtering, introduced for the first continuously operating array in southern California, estimates the shared part, the common-mode error, by stacking the residuals of a network and removes it from every station \citep{wdowinski1997}. Principal component analysis of the same network later showed that common-mode errors dominate daily solutions and are close to spatially uniform across southern California, which implies a source of very long wavelength \citep{dong2006}. Over larger areas the common mode varies: in western Europe, a filter built from the stations closest to each target removes more noise than a regional one, particularly in the vertical \citep{kreemer2021}. Part of the shared signal is physical, since atmospheric, hydrological and non-tidal ocean loading produce much of the nonlinear variation of coordinate time series \citep{li2025loading}. The station-specific part is correlated in time. Maximum-likelihood analyses describe it as white plus flicker noise in most solutions \citep{williams2004}, with random-walk, power-law and first-order Gauss--Markov components at many stations in southern California and Nevada \citep{langbein2008}; both studies find its amplitude to depend on the monument. \citet{benoist2020} fitted an exponential spatiotemporal covariance model to the stochastic variations of the IGS repro2 series and showed that accounting for spatial and temporal correlation together improves coordinates and velocities estimated from short series. These studies characterise the residual field that a virtual station must reproduce; they do not generate it where there is no station.

\subsection{Estimating fields and series away from stations}

Geodetic interpolation has mostly concerned static fields. Strain rates are interpolated from GNSS velocities with spatially varying smoothing and formal uncertainties \citep{shen2015}, and horizontal velocity fields with the Green's functions of a thin elastic sheet, which couple the two components \citep{sandwell2016}; both interpolate a field that does not change from day to day. For time series, \citet{tian2016} estimated the regional common mode with neighbours weighted by their distance and by the correlation of their long-term residuals with the target's, which requires a record at the target, and \citet{qiu2024} filled gaps in GNSS position series with an improved Gaussian process applied to the station's own record. Closest to the present task, \citet{li2023interp} predicted the series of each of five stations in North China from those of the others with gradient-boosted trees and random forests; for scattered missing days, the vertical precision improved by 45\% over cubic splines. These predictions are deterministic, and the target is an existing station with gaps. Gaussian-process regression \citep{rasmussen2006} and covariance functions for space--time data \citep{gneiting2002} provide the probabilistic counterpart, on which the kriging and Gaussian-process baselines of this study are built (Section~\ref{sec:baselines}).

\subsection{Machine learning on GNSS position series}

Deep learning on GNSS position series has so far served mostly detection and denoising. \citet{costantino2023} trained a multi-station detector of slow slip events on synthetic series made of realistic noise and modelled slip, and applied it to raw series in Cascadia from 2007 to 2022, where it retrieved 87.5\% of previously catalogued events. Graph neural networks denoise the series of irregular networks whose stations come and go: trained on the same kind of synthetic data, SSEdenoiser extracts slow slip from raw Cascadia series \citep{costantino2024}, and \citet{bachelot2025} reconstructed each station's noise from its eight nearest stations beyond 40~km, reducing the common-mode noise of the horizontal components by more than 70\% and 30\% in two processing solutions. In denoising, the part of a station that its neighbours predict is removed as noise and the remainder is kept as signal; here, that part is the starting point, and the remainder is what has to be generated. A foundation model pretrained on the series of more than 17,000 stations supports forecasting and the localisation of earthquake steps, but processes each station's window on its own \citep{teutschmann2026gnssfm}. Training on synthetic transients is thus established for detection and denoising; this study uses it to teach a generator to pass signal on (Section~\ref{sec:model}) and tests the transfer on fault-slip events and on a real earthquake.

\subsection{Kriging, imputation and diffusion models}

Machine learning has addressed the same task for other sensor networks. Inductive kriging estimates series at unobserved nodes of a sensor graph with graph neural networks trained by masking observed nodes \citep{wu2021ignnk}, and \citet{xu2025kits} reduced the gap between the sparse graph seen in training and the denser one met in inference by adding virtual nodes to the training graph. Graph neural networks also impute missing values in sensor networks \citep{cini2022grin}, including very sparse observations \citep{marisca2022spin}. These predictions are deterministic. Diffusion models generate samples by reversing a gradual noising process \citep{ho2020ddpm,karras2022edm}, and conditioning, optionally with guidance \citep{ho2022cfg}, turns them into predictors. CSDI conditions a score-based diffusion model on the observed values of a multivariate time series for probabilistic imputation and evaluates it with the CRPS \citep{tashiro2021csdi}; PriSTI adds conditional features built from spatiotemporal dependencies and geographic relationships, for spatiotemporal imputation including sensor failure \citep{liu2023pristi}; and USTD treats forecasting and kriging as conditional generation, with a deterministic encoder and diffusion decoders, on traffic and air-quality networks of 36 to 358 sensors \citep{hu2024ustd}. The model of this study belongs to this family. It differs in what geodesy requires: generation at arbitrary coordinates rather than at the nodes of a fixed graph, sparse and irregular network geometries learned by thinning a dense network, evaluation of realism (spectra, inter-station correlation, calibration across stations of different noise levels) and of the transfer of physical transients, and transfer to another region without retraining.

\subsection{Judging generated series}

A scoring rule is proper if a forecaster obtains the best expected score by reporting its true predictive distribution \citep{gneiting2007}; the continuous ranked probability score (CRPS) and its multivariate generalisation, the energy score, are proper scores for real-valued outcomes. For a finite ensemble, the usual CRPS favours ensembles drawn from too narrow a distribution, and a fair version corrects this \citep{ferro2014}. The coverage of an ensemble's intervals depends on how sample quantiles are defined \citep{hyndman1996}. Classifier two-sample tests ask whether a classifier trained to separate real from generated samples does better than chance \citep{lopezpaz2017}.

\subsection{This study}

To our knowledge, no previous study generates GNSS position series probabilistically at locations without a station from the surrounding network. Probabilistic kriging with diffusion models exists for traffic and air-quality sensors \citep{hu2024ustd}; GNSS series have been predicted from neighbours deterministically \citep{li2023interp} or for denoising \citep{bachelot2025}; and synthetic transients have trained detectors and denoisers \citep{costantino2023,costantino2024}. This study combines these elements for virtual GNSS stations and adds what they do not test: calibration across stations of different noise levels, the transfer of fault-slip events as a function of network spacing, a real earthquake, and transfer between regions.

\section{Data}
\label{sec:data}

\subsection{California}

We use the daily IGS20 solutions of NGL for 1,251 stations in and around California, from 2008-01-01 to 2026-09-26 (6,844 days), processed as in the preceding benchmark; rebuilt with this study's code from the same files, the data are identical to the benchmark's. For each station and component, a constant, a rate, annual and semi-annual terms and steps are fitted to the years 2008--2019; the residuals of this fit are the data. Equipment steps are estimated and removed; earthquake offsets are kept, since they are spatially coherent signal that a virtual station should reproduce. Outliers more than five robust standard deviations from a 61-day running median are removed, as are stations with excessive scatter (groundwater subsidence, unstable monuments) and the less complete member of any pair of stations closer than 1~km.

The years 2008--2019 are used for training, 2020--2021 for validation and 2022 onward for testing. Of the stations, 173 are held out (Figure~\ref{fig:network}). Held-out stations are excluded from training entirely: they are never a training target, never a neighbour, and never part of the input network when another station is the target. They serve only as evaluation targets. Days around the 2019 Ridgecrest earthquakes are excluded from training.

\begin{figure}[t]
\centering
\maybefigure{network_map}{0.55\linewidth}
\caption{The California network: 1,078 stations used for training and as inputs, and 173 held-out stations used only as evaluation targets.}
\label{fig:network}
\end{figure}

The residuals are heavy-tailed. Their robust scale over the training data ($1.4826$ times the median absolute deviation) is 2.12, 2.24 and 5.60~mm for east, north and up, whereas their standard deviations are 9.1, 15.3 and 10.3~mm, dominated by the coseismic offsets of the 2010 El Mayor--Cucapah earthquake, which reach about 200~mm. All series are divided by the robust scale.

\subsection{T\"urkiye}

NGL provides 233 stations with at least a year of data in the box 35.8--42.2$^\circ$N, 25.6--44.9$^\circ$E, which contains T\"urkiye and parts of its neighbours: the eastern Aegean islands, Bulgaria, Georgia, Armenia, northern Iraq and Iran. They are processed with the same steps as California's, except for the fit window. Most of the series come from T\"urkiye's national continuous network, which NGL holds densely only for 2009--2015: about 140 stations record on a typical day of those years and about 17 from 2016 on (Table~\ref{tab:turkiye}). Fitting the trajectory to 2019, as in California, would leave only the sparse remainder outside the fit. The trajectory is therefore fitted to the days up to 2012, and the main scores use 2013--2015, outside the fit and with the dense network. Of the 233 stations, 207 are kept (3 have too little data for a fit, 5 fail the scatter limits and 18 are the less complete member of a pair closer than 1~km), 176 of them in T\"urkiye according to the 1:10m country polygons of \citet{naturalearth}, with two islets in the Sea of Marmara that the polygons omit assigned to the nearest country. Every station is a target, and scores are also reported for those in T\"urkiye. The T\"urkiye data are used only for testing.

\begin{table}[t]
\centering
\caption{The T\"urkiye data: station-days, the median number of stations recording on a day, and the number of stations with at least 90 days, in each scored period. Sparse: outside the Kahramanmara\c{s} window, which runs from 30 days before to 365 days after the first mainshock. In-fit: the same stations with the trajectory fitted to 2019.}
\label{tab:turkiye}
\small
\begin{tabular}{lrrrr}
\toprule
Period & Station-days & Per day & Stations & In T\"urkiye \\
\midrule
Main: 2013--2015 & 140,138 & 141 & 167 & 144 \\
Sparse: 2016 on & 82,194 & 17 & 139 & 110 \\
Kahramanmara\c{s} & 21,750 & 29 & 108 & 92 \\
In-fit: 2009--2015 & 331,960 & 138 & 172 & 148 \\
\bottomrule
\end{tabular}
\end{table}

Three further sets of days are studied separately. From 2016 on, outside the Kahramanmara\c{s} window, the network is sparse. The period from 30 days before to 365 days after the 2023-02-06 Kahramanmara\c{s} earthquakes ($M_w$~7.8 and 7.6), whose coseismic offsets reach metres, is a case study. And the same stations with the trajectory fitted to 2019 give 2009--2015 inside the fit years, to measure what scoring inside a fit changes.

\subsection{Removal of each window's mean and trend}
\label{sec:detrend}

Residuals in the validation and test years differ systematically from those in the training years: because the trajectory was fitted to the training years, later residuals contain its extrapolation error, an offset and a drift that grow with time. Table~\ref{tab:drift} quantifies this for the held-out stations. Over 256-day windows, the slow part of the series, its 61-day running median, is about twice as large in the validation years as in the training years. After each window's least-squares mean and linear trend are removed, the two periods are indistinguishable; removing the mean alone is not sufficient for the vertical component.

\begin{table}[t]
\centering
\caption{Median root-mean-square of the slow part per 256-day window for the held-out stations, in units of each station's robust scale.}
\label{tab:drift}
\small
\begin{tabular}{lcccccc}
\toprule
& \multicolumn{3}{c}{Training years} & \multicolumn{3}{c}{Validation years} \\
\cmidrule(lr){2-4}\cmidrule(lr){5-7}
Removed from each window & E & N & U & E & N & U \\
\midrule
Nothing & 0.57 & 0.67 & 0.52 & 1.07 & 1.15 & 0.77 \\
Mean & 0.35 & 0.30 & 0.28 & 0.24 & 0.35 & 0.42 \\
Mean and trend & 0.29 & 0.21 & 0.21 & 0.18 & 0.25 & 0.20 \\
\bottomrule
\end{tabular}
\end{table}

The offset and trend are properties of the trajectory model, which is outside the scope of this work. We therefore remove from every series in a window, target and neighbours alike, its own least-squares mean and trend over its available days, in training as in evaluation, and compare generated samples with the truth after the same removal.

\section{Method}
\label{sec:method}

\subsection{Problem}

Let $\mathbf{x} \in \mathbb{R}^{L \times 3}$ be the residual series of a target location over a window of $L = 256$ days, in east, north and up. The conditioning information $\mathcal{C}$ consists of the $K = 16$ nearest stations of a network within 300~km of the target. For each neighbour $k$ it holds its series $\mathbf{y}_k \in \mathbb{R}^{L \times 3}$, its availability mask, its formal standard errors, and its east and north offset from the target in kilometres. Missing days are set to zero and flagged; unused neighbour slots are padded and masked. A neighbour counts only if at least a quarter of its window is available, so that no slot is spent on a station that is not recording. The target's absolute coordinates are not part of $\mathcal{C}$, so that a model trained in California can be applied elsewhere. The task is to draw samples from $p(\mathbf{x} \mid \mathcal{C})$.

\subsection{Training samples and network thinning}
\label{sec:thinning}

A training sample is a target station, a window lying entirely in the training years in which at least half of the target's days are available, and a network. Every admissible (station, window) pair is equally likely. The target's missing days are excluded from the loss.

Because the California network is far denser than the networks the model is meant for, each training network is thinned. In 80\% of samples the network is thinned to a median nearest-neighbour distance drawn log-uniformly from 10 to 100~km, and stations closer to the target than a distance drawn log-uniformly from 1 to 100~km are removed. The second rule prevents the model from relying on a nearby station and matches the intended use, a virtual station placed where there is none. In the remaining 20\% the full network is used with every neighbour, as in full-density evaluation: even the lightest thinning removes close pairs, whereas 13\% of the held-out stations have a training station within 5~km.

We thin by greedy minimum-separation thinning. Stations are visited in random order and kept if no kept station lies within a separation $s$, and $s$ is chosen by bisection to reach the required median spacing. Random subsampling, the obvious alternative, preserves the clustering of the California network and cannot produce a network resembling T\"urkiye's (Table~\ref{tab:thinning}): thinned at random to T\"urkiye's 60~km median spacing, California retains only 11--23 stations, with a long tail of distances beyond 100~km, whereas minimum-separation thinning retains 110--118 stations and matches T\"urkiye's median and upper tail. It has fewer close pairs than T\"urkiye (10th percentile 55 against 29~km); whether this matters will be tested on T\"urkiye rather than tuned for. Figure~\ref{fig:thinning} shows the networks used for evaluation. Thinnings are precomputed (64 spacings, four seeds each), since one thinning takes up to 0.1~s.

\begin{table}[t]
\centering
\caption{Nearest-neighbour distances of T\"urkiye's network and of California's thinned to the same median spacing, for three thinning seeds.}
\label{tab:thinning}
\small
\begin{tabular}{lrccc}
\toprule
& & \multicolumn{3}{c}{Nearest neighbour (km)} \\
\cmidrule(lr){3-5}
Network & Stations & 10th pct. & Median & 90th pct. \\
\midrule
T\"urkiye ($\geq$ 3 years) & 174 & 29 & 60 & 77 \\
California, all & 1,078 & 3 & 9 & 23 \\
California, min.-separation & 110--118 & 53--56 & 60--61 & 72--73 \\
California, random & 11--23 & 11--28 & 62--86 & 117--176 \\
\bottomrule
\end{tabular}
\end{table}

\begin{figure}[t]
\centering
\maybefigure{thinning}{\linewidth}
\caption{The training network at the four evaluation densities (kept stations dark, removed stations light), and the distribution of nearest-neighbour distances at each.}
\label{fig:thinning}
\end{figure}

\subsection{Conditional diffusion model}
\label{sec:model}

The model follows the EDM formulation of \citet{karras2022edm}. The denoiser is $D(\mathbf{x}; \sigma) = c_\text{skip}(\sigma)\,\mathbf{x} + c_\text{out}(\sigma)\,F(c_\text{in}(\sigma)\,\mathbf{x}, \tfrac14 \ln\sigma; \mathcal{C})$ with the preconditioning of \citet{karras2022edm} and $\sigma_\text{data} = 1$, close to the standard deviation of the detrended, scaled targets (about 0.9). The network $F$ is a one-dimensional U-Net over time with 64, 128, 192 and 256 channels at resolutions $L$, $L/2$, $L/4$ and $L/8$, two residual blocks per level modulated by an embedding of the noise level, and the sine and cosine of the day of year as additional inputs. The neighbours are encoded once per sample, independently of the noise level: a convolution over each neighbour's series, availability and formal errors is followed by one residual block per U-Net resolution, shared by all neighbours and modulated by an embedding of the neighbour's offset and distance (offsets scaled by 100~km, the logarithm of the distance, and sines and cosines of the offsets at wavelengths of 10, 30, 100 and 300~km). After every residual block of the U-Net, each time step of the target attends to the $K$ neighbours at the same time step; a learned null token is always among the keys, so that padded neighbour slots and dropped conditioning are handled alike. Temporal self-attention operates at the two coarsest resolutions. The network has 11.1 million parameters.

Training draws $\ln\sigma \sim \mathcal{N}(-1.2, 1.2^2)$ and minimises the weighted denoising error over the target days with data. The neighbours of 10\% of the samples are dropped, so that the model also learns the unconditional distribution and classifier-free guidance \citep{ho2022cfg} becomes available. We train for 100,000 steps of 64 samples with AdamW, a learning rate of $2 \times 10^{-4}$ with 2,000 warm-up steps and cosine decay, gradient clipping at 1 and bfloat16 arithmetic, and sample from an exponential moving average of the weights (decay 0.999). Samples are drawn by integrating the probability-flow ordinary differential equation with Heun's method over 32 noise levels from 80 to 0.002 ($\rho = 7$); the procedure is deterministic given the initial noise, which is drawn from a seeded generator, as the paired comparisons of experiment B require. On a single RTX 3060 Ti, a training step takes 0.12~s and generating 32 samples for each of about 500 target windows takes 13.7~min.

A second model, diffusion v2, adds one optional input: the target station's own noise level, the logarithm of the robust scale of its short-period variations in the training years relative to that of a typical station, for each component, with a flag saying whether it is given. A small network maps the four numbers to a vector that is added to the noise-level embedding of every block; its last layer starts at zero, so that training starts from the first model. The level is hidden from 30\% of the training samples, so that one model generates with or without it. Section~\ref{sec:noise} motivates the input. In use, it lets the user choose how noisy a virtual station should be; in evaluation, the held-out station's training-years level stands in for that choice, which uses information about the real station, and results with and without it are reported separately. Diffusion v2 is otherwise trained exactly as the first model.

A third model, diffusion v3, is trained with synthetic transients planted in the training samples, so that it learns to pass on localised signals that the training years rarely contain. With probability 0.8 a sample is drawn with one transient, added to the target and to every neighbour before the removal of each window's mean and trend. Its displacement field has a Gaussian footprint of radius 10--300~km (log-uniform), centred within 1.5 radii of the target, so that the target lies inside it, at its edge or outside it. In half of the transients the displacement has one direction everywhere (random azimuth, vertical share up to $\pm 0.5$); in the other half it is radial about the centre, with the horizontal displacement largest at one radius and the vertical at the centre, so that neighbours on opposite sides move in opposite directions. In time, 70\% are raised-cosine rises over 5--120~days with amplitudes of 1--30~mm (log-uniform), of which 30\% return to zero over one to three times the rise, and 30\% are steps with heavy-tailed amplitudes (Pareto, $\alpha = 1.5$, from 1~mm, capped at 50~mm); onsets lie between a quarter-window before the window and its end. A transient is kept only if at least two neighbours change by 1~mm or more over their available days. Otherwise the target could carry a transient that nothing in the input reveals, and the model would learn to add unexplained excursions. About 40\% of the drawn transients fail this test, so that 49\% of the training samples carry one. The family is deliberately simpler than the fault-slip events of experiment B, which therefore test generalisation rather than recall. Diffusion v3 is otherwise trained exactly as diffusion v2, with the same number of steps and the same seed; regress v2 is regress trained with the same transients and the station's noise level as an optional input.

The regression counterpart, \textbf{regress}, is the same network evaluated with a zero input and noise level and trained for 50,000 steps with the squared error on the target days with data. It is the conditional mean that the diffusion model would reach if it predicted a single value.

\subsection{Baselines}
\label{sec:baselines}

Six baselines define what the diffusion model has to improve on: three simple ones, and three built on a covariance model of the residual field fitted to the training years.
\begin{description}
\item[Zero] predicts zero.
\item[IDW] predicts the weighted mean of the available neighbours on each day, with weights $w_k = \exp(-d_k^2 / 2\ell^2) / \bar\sigma_k^2$, where $d_k$ is the distance to neighbour $k$ and $\bar\sigma_k$ its median formal standard error in the window. The length scale $\ell$ is selected for each density among 10, 25, 50, 100, 200 and 500~km and $\infty$ by the short-period error on validation days, with 200 training stations as targets, each predicted without itself, so that the held-out stations play no part in the choice.
\item[IDW + noise] adds to the IDW prediction white and flicker noise \citep{williams2004}, independently per component, with periodogram level $P(f) = A + B/f$. For each window, $A$ and $B$ are fitted by non-negative least squares over 1/64 to 0.5 cycles per day to the departure of each neighbour from its own IDW prediction, and the medians over the neighbours are used. Fitting both terms jointly matters, since flicker noise retains power at high frequencies, and least squares is preferred to medians of periodogram ordinates, whose exponential distribution would bias the levels low by a factor $\ln 2$. The noise is synthesised on twice the window length and truncated, so that it does not wrap around. This is the simplest generative model of a virtual station.
\item[Kriging] predicts each day by Gaussian-process regression \citep{rasmussen2006} from the neighbours with data on that day. Each component has the stationary, isotropic covariance $C(d) = s_\mathrm{cm} + s_\mathrm{r} M(d/\ell) + s_\mathrm{n} \delta$, where $M$ is the Mat\'ern correlation with $\nu = 3/2$, $s_\mathrm{cm}$ a common mode shared by the whole network, $s_\mathrm{r}$ a regional part and $s_\mathrm{n}$ the station's own variance; the common-mode term acts as the unknown mean of ordinary kriging with a fitted prior. The covariance is fitted once, to the training years of the training stations, and serves every density: the mean covariance of all station pairs within 300~km in 14 distance classes and the mean station variance, over 256-day windows detrended as the model's inputs, are fitted by weighted least squares, with $\ell$ by grid search. Station-windows with a listed earthquake step larger than one robust scale are left out. Nearly all the shared covariance is network-wide: it is constant from 25 to 300~km, and the regional part is small and confined to the first 20~km ($s_\mathrm{cm}$ = 0.26, 0.14 and 0.18 of station variances of 0.65, 0.41 and 0.80 for east, north and up, in scaled units).
\item[Kriging + noise] adds noise to the kriging prediction as IDW + noise does, fitted to the departures of the neighbours from their own kriging predictions.
\item[Space--time Gaussian process] generates the target's whole window jointly. Each component is the sum of three independent parts, each separable in space and time: the common mode, the regional part (Mat\'ern, $\ell$ from the kriging fit) and the station's own part, scaled by the square of the station's noise level relative to a typical station. The temporal covariance of each part is the sum of white noise, flicker noise and a first-order Gauss--Markov process \citep{williams2004,langbein2008}; without the Gauss--Markov term the station part falls short of the observed covariance at lags of 6--16 days. The parts are fitted to lag curves of the training years between 0 and 128 days, computed exactly for a detrended window: from pairs 100--300~km apart (common mode), pairs within 10~km (regional) and each station with itself (station part), with amplitudes by non-negative least squares and the Gauss--Markov time constant by grid search. Space--time interaction \citep{gneiting2002} would act only on the small regional part and is omitted. Samples are drawn from the target's distribution conditional on every available neighbour value in the window, with a flat prior on each neighbour's mean and trend \citep[eq.~2.42]{rasmussen2006}, which conditions on the neighbours as detrended. Covariances in the lag curves are medians over stations or station pairs within each window, averaged over windows; Section~\ref{sec:gaussian} explains why.
\end{description}
Related methods that are not reimplemented here are reviewed in Section~\ref{sec:related}.
A fourth comparison, the diffusion model's network trained with a squared-error loss, will show what the generative formulation adds beyond its own conditional mean.

\section{Evaluation}
\label{sec:eval}

\subsection{Experiment A: held-out stations}

Each of the 173 held-out stations is generated from the training stations, thinned to one of four densities: the full network, and median spacings of 25, 60 and 100~km. The 60~km density matches T\"urkiye and is the main result. Each thinned density is repeated with five thinning seeds. The validation period is tiled with 256-day windows; the last window is aligned to the end of the period and scores only the days not covered by an earlier window. All decisions use the validation period; the test period is used once, for the final report.

Errors are expressed in units of each station's robust scale over the training years, every score is computed per station, and stations are summarised by the median, following the preceding benchmark. Days within 30 days of an earthquake of magnitude 6 or more that NGL lists as possibly offsetting the station are excluded from the short- and long-period scores, and 19 held-out stations whose scatter increased more than threefold, or whose series jumped without a listed cause, after 2019 are excluded from scoring.

\begin{description}
\item[Point accuracy.] The error of the sample mean is split into a long-period part, its 61-day running median, and a short-period part, the remainder; the normalised root-mean-square error (nRMSE) of each, and of the total, is reported.
\item[Probabilistic accuracy.] The continuous ranked probability score (CRPS) and its multivariate counterpart, the energy score \citep{gneiting2007}, are computed with the fair ensemble estimators of \citet{ferro2014}, which are unbiased for the distribution the samples are drawn from. For a point forecast the CRPS reduces to the absolute error. Coverage of the central 50\% and 90\% intervals is computed with Weibull plotting positions \citep{hyndman1996}, under which a draw from the ensemble's own distribution falls inside the 90\% interval with probability 0.9; the default positions of common software give about 0.85 for 32 members. The ratio of ensemble spread to the error of the ensemble mean is 1 for a calibrated ensemble. For the short-period scores the truth and each sample are decomposed separately. Because a median over stations emphasises the typical station, calibration is also reported pooled over all scored days of all stations, and separately for the quartiles of the held-out stations' own noise level, the robust scale of their short-period variations in the training years relative to that of a typical station. The observed series were screened for outliers when the data cube was built (Section~\ref{sec:data}); for a like-for-like comparison, generative methods are also scored after the same screen is applied to every sample, with a flagged value replaced by the sample's running median.
\item[Realism.] Single samples are compared with the truth in four ways: the ratio of their power spectral densities in six period bands from 2 to 128 days, the ratio of the robust scales of their short-period parts, the ratio of the 99th percentiles of their absolute short-period values, and their correlation with the neighbours as a function of distance.
\item[Indistinguishability.] A classifier trained to separate real from generated windows \citep{lopezpaz2017} will be added with the diffusion model.
\end{description}
Differences between two methods are tested on the same stations and networks. The difference of the median summaries has a 95\% interval from a bootstrap over stations (10,000 resamples), in which a station keeps all its thinned networks together, and the number of stations at which each method is better is tested with an exact sign test. Both treat stations as independent; nearby stations share signal, so the intervals are somewhat optimistic, while the counts of stations are robust to it.

\subsection{Experiment B: planted fault-slip events}

To test whether localised transients are transferred to a virtual station, slow-slip and afterslip events are placed on fault surfaces of the SCEC Community Fault Model \citep{plesch2007,plesch2024cfm}, and their surface displacements are computed with triangular dislocations in an elastic half-space \citep{nikkhoo2015}, with Poisson's ratio 0.25. Two templates follow published events: shallow slow slip on steep strike-slip faults, 8--25~km long, from the surface to 2--4~km depth, with 0.8--4.5~cm of slip propagating along strike at 0.4--9~km per day \citep{materna2024,vavra2024}; and afterslip on 15--50~km patches between the surface and 5--12~km depth, with 10--40~cm of slip growing as $\log(1 + t/\tau)$ over a year, $\tau$ = 17--34~days \citep{gonzalezortega2014}. Thirty events of each template are placed on random suitable faults, with onsets in the validation period such that each event and a further 60 days end within it, and each is kept only if at least one scored held-out station moves by more than 1~mm horizontally. The radius of the region in which an event's final horizontal displacement exceeds 1~mm is its footprint: 4--17~km (median 9~km) for slow slip and 29--122~km (median 53~km) for afterslip.

Each event is added to every station's real series, the held-out targets included, and every held-out target that it moves by more than 1~mm is generated twice over the windows of experiment A that overlap the event, once from the network with the event and once from the original network. Both runs use the same random numbers for each target and window, whatever else is generated with them, so that the difference $\mathbf{d}$ of the two sample means (16 samples) isolates what the method passed on from the event. Detrending is linear, so the planted signal in the detrended truth is the detrended event displacement $\mathbf{s}$ at the target. The recovered fraction is $\rho = \langle\mathbf{d}, \mathbf{s}\rangle / \langle\mathbf{s}, \mathbf{s}\rangle$ over the horizontal components and the event's days (from the onset to 60~days after its end), reported per target and event and pooled over them weighted by $\langle\mathbf{s}, \mathbf{s}\rangle$. The shape (correlation of $\mathbf{d}$ with $\mathbf{s}$), the CRPS and coverage of the planted truth by the planted samples, and the spurious signal at up to five held-out targets that the event leaves still but whose neighbours it moves are also reported. Per-window detrending removes most of a signal that grows over months: an afterslip event that moves a station by 12~mm leaves a detrended signal of 0.4~mm r.m.s., so $\rho$ measures the transfer of what survives detrending, as for the data. Experiment B uses the full network and two thinnings at each of 25, 60 and 100~km.

As a real event, held-out stations within 300~km of the 2019 Ridgecrest earthquakes are generated over a 256-day window centred on them, and the coseismic offset of the samples (mean of 10 days after minus 10 days before, 2 days either side skipped) is compared with the offset observed at the station. The window lies in the training years, but the held-out stations never entered training, and the stations near the rupture were excluded from it in the second half of 2019.

\subsection{Experiment C: T\"urkiye}

Every station is generated in turn from all the others, its 16 nearest within 300~km, at the network's own density. Nothing is tuned on T\"urkiye: the trained models and their input scaling, the length scale of IDW (California's at 60~km, T\"urkiye's spacing) and the covariances of kriging and the space--time Gaussian process all come from California, and a station's noise level is relative to California's typical station, as the models learned it. Scoring follows experiment A with three changes: the quality control that drops stations for site trouble checks the scored period's own days, so that a station is not dropped from 2013--2015 for trouble years later; stations with too short a record before the fit's end, whose trajectory is fitted to the start of their record and so overlaps the scored days, are left out; and the Kahramanmara\c{s} window is skipped. The Kahramanmara\c{s} sequence is analysed separately: the coseismic offsets and the first year of postseismic deformation of virtual stations near the rupture are compared with the observed ones, as a function of distance from the rupture and from the nearest station.

\section{Results}
\label{sec:results}

\subsection{Baselines on held-out stations}

Table~\ref{tab:a} and Figure~\ref{fig:a} summarise experiment A on the validation period for the three baselines, the regression network and the diffusion model. Of the 173 held-out stations, 154 are scored; the others fail the scoring quality control or have fewer than 20 scored days. We first describe the baselines, which set the reference, and then the trained models.

\begin{table}[t]
\centering
\caption{Experiment A on the validation period (2020--2021): median over 154 held-out stations, mean over five thinning seeds for the thinned densities. Subscripts $f$ and $t$ denote the fast part and the total; Cov$_{90,f}$ is the coverage of the 90\% interval of the fast part, Spread/err.\ the ratio of ensemble spread to the error of the ensemble mean, and log$_2$ scale the ratio of the short-period scale of single samples to that of the truth. Errors are in units of each station's robust scale; lower is better for nRMSE and CRPS, and 0.9, 1 and 0 are ideal for coverage, spread/error and log$_2$ scale. Coverage, spread and scale are not defined for point forecasts.}
\label{tab:a}
\footnotesize
\setlength{\tabcolsep}{3.5pt}
\maybetable{experiment_a_table}
\end{table}

\begin{figure}[t]
\centering
\maybefigure{scores_by_density}{\linewidth}
\caption{Scores against network density. Lines are means over five thinning seeds and bands their range; dotted lines mark the ideal value.}
\label{fig:a}
\end{figure}

\paragraph{Point accuracy depends little on density} Predicting a held-out station from its neighbours removes much of its short-period error: IDW reduces the fast nRMSE from 0.737 for predicting zero to 0.527 with the full network. The reduction hardly depends on density. At T\"urkiye's density of 60~km the fast nRMSE is 0.532, and even at 100~km it is 0.547. The length scales selected on validation are long (200~km to infinite), so all neighbours are weighted nearly equally. Both observations indicate that the short-period variability a station shares with its neighbours is dominated by common-mode error coherent over hundreds of kilometres \citep{wdowinski1997,dong2006}, which a sparse network samples almost as well as a dense one. Consequently, the short-period point scores of experiment A are only weakly sensitive to network density, and what densification contributes beyond the common mode has to be measured on spatially localised signals, the purpose of experiment B. After each window's mean and trend are removed, the slow part contributes little error (nRMSE 0.13 for IDW, 0.25 for zero).

\paragraph{A noise model makes the mean a calibrated generator} Adding white and flicker noise to the IDW mean lowers the fast CRPS from 0.393 to 0.290 at full density and from 0.399 to 0.293 at 60~km, a reduction of 26\%. The mean of its 32 samples is slightly less accurate than IDW itself (fast nRMSE 0.536 against 0.527), since the mean of a finite ensemble retains part of the added noise. The samples are close to calibrated: the 90\% intervals cover 90--91\% of the fast part and the 50\% intervals 53--54\%; the ratio of spread to error of the total is 0.91; and the short-period scale of single samples exceeds that of the truth by 5--8\% ($\log_2$ ratio 0.07--0.11). Their spectra lie within $\pm 0.06$ in $\log_{10}$, about $\pm 15\%$, of the truth in every band from 2 to 128 days, with a consistent excess of 12--15\% at periods of 2--4 days in all components (Figure~\ref{fig:spectra}).

\begin{figure}[t]
\centering
\maybefigure{spectra_60km}{\linewidth}
\caption{Power of single samples relative to the truth by period band, at 60~km density (thinning seed 0); 0 is a perfect match.}
\label{fig:spectra}
\end{figure}

\paragraph{The spatial structure is not reproduced} The correlation of a real station's short-period variations with those of its neighbours decreases with distance, from 0.50 within 10~km to 0.45--0.48 beyond (Figure~\ref{fig:corr}). The samples of IDW + noise correlate with their neighbours at 0.48--0.53 at all distances: more strongly than real stations do with distant neighbours, by 0.035--0.048 at full density, and without the decrease. A near-uniform weighted mean with independent noise reproduces the average amount of shared signal but not how it varies in space. Reproducing this structure is a specific requirement for the diffusion model. Figure~\ref{fig:example} shows one held-out station with its samples.

\begin{figure}[t]
\centering
\maybefigure{correlation_60km}{0.6\linewidth}
\caption{Correlation of the short-period parts of target and neighbour against their distance, for the truth and for the samples of the generative methods, at 60~km density (thinning seed 0). Within 10~km there are only 24 pairs at this density.}
\label{fig:corr}
\end{figure}

\begin{figure}[t]
\centering
\maybefigure{example_idw_noise}{\linewidth}
\caption{Held-out station ARTE in the first validation window, generated by IDW + noise from the network thinned to 60~km: eight of 32 samples, their mean, and the truth, after removal of the window's mean and trend.}
\label{fig:example}
\end{figure}

\subsection{The diffusion model and its regression counterpart}

\begin{figure}[t]
\centering
\maybefigure{training_curves}{\linewidth}
\caption{Training curves: the training loss every 100 steps (faint) and its 1,000-step running mean, and the loss on 1,024 fixed validation samples of held-out stations, evaluated with the moving-average weights. The two losses are on different scales.}
\label{fig:training}
\end{figure}

\paragraph{Training} The validation loss of the diffusion model fell from 0.666 after 2,000 steps to 0.594 after 20,000 and 0.582 at the end of training, still decreasing slowly as the learning rate decayed (Figure~\ref{fig:training}). The regression network reached its lowest validation error, 0.476, within 10,000 steps and changed by less than 0.2\% thereafter, while its training error kept falling; its conditional mean is learned quickly and the remaining steps fit the training stations.

\paragraph{Accuracy} The regression network is the most accurate point predictor at every density, with a fast nRMSE 1.5--2\% below that of IDW (0.519 against 0.527 with the full network, 0.524 against 0.532 at 60~km). The improvement over a tuned weighted mean is small, consistent with the conclusion of Section~\ref{sec:results} that the predictable short-period signal is dominated by a common mode that a near-uniform average already captures. The mean of 32 diffusion samples is less accurate (0.534 and 0.536), as that of IDW + noise is, because a finite ensemble mean retains part of the sampled variability.

\paragraph{Probabilistic accuracy} The diffusion model has the lowest short-period CRPS at every density: 0.279, 0.281, 0.285 and 0.291 for the full network and for 25, 60 and 100~km, against 0.290, 0.291, 0.293 and 0.301 for IDW + noise, a reduction of 3--4\%. Its energy score over the three components is also the lowest (0.593 against 0.613 at 60~km). At T\"urkiye's density, the target set by the baselines, a short-period CRPS below 0.29, is met (0.285, range 0.284--0.286 over five thinnings).

\paragraph{Calibration} In the short-period part the diffusion samples are close to calibrated in the bulk: the 50\% intervals cover 50--51\% of the truth, and the robust scale of single samples matches that of the truth to within 1--2\% ($\log_2$ ratio 0.004--0.026, against 0.08--0.11 for IDW + noise). At the median station, however, the 90\% intervals cover 92.5--93\%, and over the whole series the ensemble spread exceeds the error of the ensemble mean by 21--26\%.

\paragraph{Spatial structure} Unlike IDW + noise, the diffusion samples reproduce the decrease of inter-station correlation with distance (Figure~\ref{fig:corr}). At 60~km density their correlation with neighbours falls from 0.507 within 10~km to 0.44--0.46 beyond 25~km, against 0.512 and 0.45--0.47 for the truth; at every distance and density the difference from the truth is at most 0.023 and mostly below 0.015, whereas IDW + noise exceeds the truth by 0.03--0.05 beyond 10~km. The model has learned how the shared signal varies with distance, the requirement identified in Section~\ref{sec:results}.

\paragraph{Excess variability} The diffusion samples carry 3--16\% more power than the truth in most period bands (Figure~\ref{fig:spectra}), and 22\% more in the vertical at 64--128 days, and the 99th percentile of their absolute short-period values exceeds that of the truth by 21\% at 60~km density ($\log_2$ ratio 0.27). Figure~\ref{fig:example_diffusion} shows the form the excess takes at one station: some samples drift slowly by up to 10~mm over months, or jump, without any counterpart in the truth or in the neighbours. Two explanations suggested themselves, and both turned out to be minor. First, the observed series were screened for outliers and the samples were not. Applied to the samples, the screen flags 0.07\% of their values (against 0.5\% of the observed station-days when the cube was built) and leaves every score nearly unchanged (Table~\ref{tab:variants}): the excess is not made of isolated spikes. Second, the training residuals retain the coseismic offsets of earthquakes in 2008--2019. Windows in which the target moves by more than one robust scale at a listed earthquake are 6\% of the training samples but carry 22\% of their variance, so the model might have learned that a station occasionally moves on its own. If so, it would be too wide on the quiet validation period but not on the training distribution. It is too wide on both. Generating eight samples for 1,536 windows drawn as in training and for 1,536 validation windows, the median ratio of spread to error per window is 1.14 for training windows without an earthquake step, 1.12 for those with one and 1.11 for validation windows, and the samples' long-period parts exceed the truth's by 8--10\% in all three groups.

\paragraph{The model cannot know how noisy a station is} Sorting the validation windows by the noise of the truth exposes the cause. The 99th percentile of the truth's absolute short-period values rises from 1.30 in the quietest quarter of windows to 2.20 in the noisiest; that of the samples rises only from 1.60 to 1.86, and the ratio of spread to error falls from 1.28 to 0.84. The training windows show the same pattern. The amplitude of a station's own noise depends on its monument, antenna and environment, which its neighbours do not reveal, and the model hedges over the stations that could stand at a location: its ensembles are too wide where the station is quiet and too narrow where it is noisy. Across the whole validation grid (Figure~\ref{fig:calibration}), the ratio of spread to error of the diffusion model falls from 1.43 for the quietest quarter of held-out stations, ranked by their noise in the training years, to 0.97 for the noisiest, and the coverage of its 90\% short-period intervals from 0.96 to 0.87. Pooled over all stations, the short-period intervals are calibrated (coverage 0.909--0.912 at every density) and the ratio of spread to error of the total is 1.11--1.18, so a residual excess of 10--15\% remains at long periods. A median over stations, in which the typical station is quieter than the average, makes the ensembles look wider than they are. IDW + noise shows the same gradient, from 1.08 to 0.67, at a lower level: its noise model is fitted to the neighbours and cannot know the target's noise either.

\begin{figure}[t]
\centering
\maybefigure{calibration_by_noise_60km}{\linewidth}
\caption{Calibration of the generative methods by the held-out station's own noise level at 60~km density: quartiles of the robust scale of the station's short-period variations in the training years, relative to a typical station (median of each quartile below its label), and all stations pooled (open markers). Lines are means over five thinning seeds and bands their range; dotted lines mark the ideal value.}
\label{fig:calibration}
\end{figure}

\paragraph{Classifier-free guidance} Sampling with a guidance weight $w > 1$ amplifies the difference between the conditional and the unconditional prediction. At $w = 1.2$ the ratio of spread to error of the total reaches 1.02, but the short-period intervals become too narrow (coverage 0.85), the CRPS rises from 0.284 to 0.289 and the samples correlate with their neighbours by 0.07 more than real stations do; at $w = 1.5$ every score is worse than without guidance (Table~\ref{tab:variants}). Guidance strengthens the shared signal that the model attributes to the neighbours, whereas the excess lies in the station's own variability; it is not used further.

\begin{table}[t]
\centering
\caption{Sampling variants at 60~km density (thinning seed 0): the guidance weight $w$, and scores with and without the outlier screen applied to the samples. S/E is the ratio of spread to error of the total at the median station and S/E$_\mathrm{pool}$ pooled over stations; Tail the $\log_2$ ratio of the 99th percentiles of absolute short-period values of samples and truth; $|\Delta r|$ the mean absolute error of the neighbour correlation over distance classes; Scr.\ the share of sample values flagged by the screen.}
\label{tab:variants}
\footnotesize
\setlength{\tabcolsep}{3.5pt}
\maybetable{variants_table}
\end{table}

\begin{figure}[t]
\centering
\maybefigure{example_diffusion}{\linewidth}
\caption{The same window as Figure~\ref{fig:example}, generated by the diffusion model: eight of 32 samples, their mean, and the truth. Two samples drift in the north component by up to 10~mm over several months and one jumps in the vertical, without counterpart in the truth.}
\label{fig:example_diffusion}
\end{figure}

\subsection{Comparison with Gaussian generators}
\label{sec:gaussian}

\paragraph{Kriging adds nothing to the weighted mean} Kriging matches IDW at every density (fast nRMSE 0.529, 0.527, 0.533 and 0.548 for the full network and 25, 60 and 100~km, against 0.527, 0.527, 0.532 and 0.547), and kriging + noise matches IDW + noise (fast CRPS 0.291, 0.291, 0.294 and 0.302, against 0.290, 0.291, 0.293 and 0.301). This is what the fitted covariance implies: with a shared signal that is almost constant over 300~km, the optimal linear weights are nearly uniform, as the long length scales selected for IDW already indicated, and the slight shrinkage of kriging towards zero costs 0.001--0.002. The diffusion model's lead is therefore not a lead over a naive average only; the best linear predictor under a covariance fitted to the data does no better than the average.

\paragraph{The decrease of correlation with distance is not specific to the diffusion model} Kriging + noise reproduces it as well as the diffusion model does (Figure~\ref{fig:corr}; mean absolute error over distance classes 0.007--0.017 across densities, against 0.008--0.012 for the diffusion model and 0.033--0.052 for IDW + noise): the small regional covariance gives near neighbours more weight. The failure of IDW + noise in this respect follows from its near-uniform weights rather than from a limitation of Gaussian generators.

\begin{table}[t]
\centering
\caption{Generative methods on one network (60~km density, thinning seed 0), unscreened. Energy is the energy score; S/E the ratio of spread to error at the median station and S/E$_\mathrm{pool}$ pooled over stations; Scale and Tail the $\log_2$ ratios of the robust scale and of the 99th percentile of single samples' short-period values to the truth's; PSD$_{64\text{--}128}$ the $\log_{10}$ ratio of power at periods of 64--128 days; $|\Delta r|$ the mean absolute error of the neighbour correlation over distance classes.}
\label{tab:generators}
\footnotesize
\setlength{\tabcolsep}{3pt}
\maybetable{generators_table}
\end{table}

\paragraph{A space--time Gaussian process fitted to typical covariances} The first fit of the space--time Gaussian process used means of the covariances over station-windows. Its samples were 25\% too large at short periods, 2.6 times too strong at 64--128 days and too weakly correlated with their neighbours, and its CRPS (0.306) was worse than that of IDW + noise (Table~\ref{tab:generators}). The cause lies in the heavy tails of the residuals: in the training years the mean variance of a 256-day window is 1.5--1.7 times the median, and that of its long-period part 5--6 times, so that a Gaussian model fitted to means describes a disturbed station rather than a typical one; it is also inconsistent with station noise levels, which are relative to the median station. Fitted instead to medians within each window, the Gaussian process matches IDW + noise in CRPS (0.293) and in energy score (0.612), has the best calibration at the median station (ratio of spread to error 0.99, coverage of 90\% intervals 0.92), reproduces the spectra within 0.02--0.05 in $\log_{10}$ and the decrease of correlation with distance (mean error 0.009). The change of statistic was prompted by the first validation result, although both fits use the training years alone; both are therefore reported. Over all networks the Gaussian process fitted to medians scores a fast CRPS of 0.289, 0.291, 0.293 and 0.299 for the full network and 25, 60 and 100~km, within 0.002 of IDW + noise, with 90\% coverage of 0.92 and a ratio of spread to error of 0.99--1.00 at the median station at every density.

\begin{table}[t]
\centering
\caption{Paired comparisons of fast CRPS on the same held-out stations and networks: relative difference (\%) of the median summaries with its 95\% bootstrap interval, and the number of the 154 scored stations at which the first method is better. Negative values favour the first method. Thinned densities pool five networks; ``given'': with the held-out station's noise level.}
\label{tab:significance}
\scriptsize
\maybetable{significance_table}
\end{table}

\paragraph{What the diffusion model adds} The diffusion model improves on the best Gaussian generator at every density, by 2.6--3.4\% in short-period CRPS and 2.0--3.1\% in energy score, and it is better at 136--145 of the 154 scored stations (Table~\ref{tab:significance}; every 95\% interval excludes zero, and every sign test gives $p < 10^{-19}$), with the same agreement in neighbour correlation (0.008--0.012 against 0.009--0.015). Its ensembles are wider than the Gaussian process's at the median station (ratio of spread to error 1.26 against 0.99) but closer to calibrated pooled over stations (1.16 against 0.83), and its samples have heavier tails than the truth (Tail 0.27) where the Gaussian samples have lighter ones. The gain from a learned, non-Gaussian generator is modest for held-out stations under ordinary conditions but consistent across stations and densities; whether it is larger for localised transients, which a stationary covariance cannot represent, is the question of experiment B (Section~\ref{sec:resultsb}).

\subsection{Specifying the station's noise level}
\label{sec:noise}

\begin{figure}[t]
\centering
\maybefigure{calibration_noise_models_60km}{\linewidth}
\caption{Calibration by the held-out station's noise quartile at 60~km density, as Figure~\ref{fig:calibration}, for the first diffusion model, diffusion v2 with and without the station's noise level, and the space--time Gaussian process with and without it; means over five thinnings.}
\label{fig:noisemodels}
\end{figure}

\begin{table}[t]
\centering
\caption{The diffusion models and the Gaussian process with and without the station's noise level; means over thinning seeds. S/E$_{Q1}$ and S/E$_{Q4}$ are the ratios of spread to error for the quietest and noisiest quarters of held-out stations.}
\label{tab:noisemodels}
\scriptsize
\setlength{\tabcolsep}{2.5pt}
\maybetable{noise_models_table}
\end{table}

Given the held-out station's noise level, diffusion v2 is calibrated across the quartiles of station noise: the ratio of spread to error lies between 0.96 and 1.14 from the quietest to the noisiest quarter at both densities, instead of falling from 1.43--1.46 to 0.97, and the coverage of its 90\% short-period intervals is 0.90--0.92 in every quarter (Figure~\ref{fig:noisemodels}, Table~\ref{tab:noisemodels}). The same input removes most of the excess variability of the first model: the 99th percentile of the samples' absolute short-period values exceeds the truth's by 3\% instead of 21\% ($\log_2$ ratio 0.044 against 0.27 at 60~km), and the excess power at 64--128 days disappears ($\log_{10}$ ratio 0.001 against 0.052). The excess of the first model was therefore almost entirely the hedge over the station's unknown noise. The fast CRPS improves only slightly (0.278 against 0.279 with the full network, 0.283 against 0.285 at 60~km), but the mean of the samples becomes more accurate (fast nRMSE 0.527 and 0.529 against 0.534 and 0.536), since the samples no longer wander more than the station does.

Without the noise level, diffusion v2 scores as the first model (fast CRPS 0.279 and 0.284; paired differences of $-0.2$ and $-0.4\%$ with intervals spanning zero), so hiding the input from part of the training costs nothing; it is nonetheless better calibrated than the first model (ratio of spread to error 1.12 against 1.23 at 60~km, 0.98 against 1.14 pooled), as if learning the input had helped it separate station noise from shared signal. The Gaussian process gains little from the same information: scaling the target's station part by its noise level leaves its CRPS unchanged (0.290 and 0.292) and only partly flattens its calibration across quartiles (with the full network, from 1.09 to 0.80 instead of 1.17 to 0.71), because in a Gaussian model the noise level changes the width of the ensemble but not the shape of its tails or spectrum. Given the same information, diffusion v2 improves on the Gaussian process by 3.9\% in CRPS with the full network and 3.1\% at 60~km (Table~\ref{tab:significance}).

\subsection{Experiment B: planted events and Ridgecrest}
\label{sec:resultsb}

\begin{figure}[t]
\centering
\maybefigure{b_rho_by_ratio}{0.65\linewidth}
\caption{Recovered fraction of planted events (median over target--event pairs) against the event footprint divided by the distance from the target to its nearest network station, pooled over the full network and the six thinnings; below, the number of target--event pairs per class.}
\label{fig:brho}
\end{figure}

\begin{table}[t]
\centering
\caption{Recovered fraction of planted events, pooled over target--event pairs weighted by the energy of the detrended planted signal. Thinned densities pool two thinnings.}
\label{tab:brho}
\footnotesize
\maybetable{b_rho_table}
\end{table}

\paragraph{Only events larger than the gaps between stations reach a virtual station} No method passes on a planted event whose footprint is smaller than about twice the distance from the target to its nearest network station (Figure~\ref{fig:brho}): if no neighbour moves, nothing can be transferred. Shallow slow slip, with footprints of about 9~km, is therefore invisible at every density (pooled $\rho \le 0.03$; Table~\ref{tab:brho}), and so is afterslip once the network is thinned to 60~km or more ($\rho \le 0.03$). The transfer of localised signals is where densification matters, in contrast to the common-mode noise of experiment A.

\paragraph{The Gaussian process transfers the most} Over all networks, where the event is more than four times larger than the gap to the nearest station, the space--time Gaussian process recovers a median of 57\% of the planted signal, against 44\% for diffusion v2, 42\% for regress and 27--34\% for IDW, IDW + noise and kriging + noise. Its margin is widest on the full network (91\% against 72--77\% for the others), where diffusion v2 passes on no more than the near-uniform averages; on the thinned networks diffusion v2 passes on more than the averages (47\% against 38--41\% at 25~km) but less than the Gaussian process (66\%). Pooled over afterslip events on the full network, it recovers 31\% of the planted signal energy, against 26\% for regress and 17--19\% for the diffusion model, IDW and kriging, and the shape of what it passes on is the closest to the event (median correlation 0.98 against 0.89 for the diffusion model). Conditioning on the whole window, with a regional component correlated over months, suits a slow, regionally coherent signal. The diffusion model, trained on years with few localised transients, passes on less; its non-Gaussian flexibility does not by itself translate into transferring local signals. During the events it keeps its probabilistic lead (CRPS of the planted truth 0.300 against 0.306 for the Gaussian process with the full network, coverage of 90\% intervals 0.83 for both) and adds as little spurious signal at stations outside the event as the other methods (median over all networks 1.1--1.5\% of the neighbours' peak displacement; 2.4\% for regress).

\begin{table}[t]
\centering
\caption{The 2019 Ridgecrest earthquakes at 13 held-out stations within 300~km whose observed horizontal offset exceeds 5~mm: median recovered fraction of the observed offset, median error of the generated offset, and the share of offset components inside the 90\% range of the samples.}
\label{tab:ridgecrest}
\footnotesize
\setlength{\tabcolsep}{2pt}
\maybetable{b_ridgecrest_table}
\end{table}

\paragraph{A real earthquake} Of 39 held-out stations within 300~km of the Ridgecrest earthquakes with data around them, 13 moved by more than 5~mm (Table~\ref{tab:ridgecrest}). With the full network the averaging methods reproduce their offsets (recovered fraction 1.0--1.1, median error 2.7--3.7~mm), the space--time Gaussian process most closely, while the diffusion models recover about half to two thirds (0.53--0.64, errors 5--7~mm). Thinned to 60~km, every method recovers little (0.09--0.24); the first diffusion model, regress and the Gaussian process recover the most (0.20--0.24), and the Gaussian process has the smallest error (9.1~mm against 9.4--9.8~mm for the averages and 10.9--12.2~mm for the diffusion models). The diffusion models are the only methods whose ranges contain the observed offsets at a useful rate (54--77\% of components inside the 90\% range, against 8--19\% for the noise-adding baselines and the Gaussian process), so their ensembles acknowledge an uncertainty that the Gaussian generators, built for ordinary days, understate. Near-field offsets of this size were excluded from training, and the diffusion models appear to pass large steps on only in part.

\subsection{Training with planted transients}
\label{sec:transients}

\begin{table}[t]
\centering
\caption{Experiment A for the models trained with and without planted transients; means over thinning seeds. Columns as in Table~\ref{tab:variants}; ES is the energy score and Cov.$_{90,f}$ the coverage of the fast part's 90\% intervals at the median station.}
\label{tab:transients}
\scriptsize
\setlength{\tabcolsep}{2.5pt}
\maybetable{transients_table}
\end{table}

\paragraph{Ordinary days are unaffected} The validation loss of diffusion v3 ended at 0.581, against 0.582 for diffusion v2, and its loss on the same validation samples with planted transients at 0.588. In experiment A it scores as diffusion v2 (Table~\ref{tab:transients}): the paired differences in fast CRPS lie between $-0.4$ and $+0.3\%$, with 95\% intervals that include zero, with and without the station's noise level, at full density and at 60~km. Its energy score at 60~km is 0.5--0.8\% higher; the intervals include zero, but diffusion v2 is better at 94--95 of the 154 stations ($p < 0.01$), so the planted transients cost a small amount of multivariate accuracy. Given the noise level, diffusion v3 is calibrated across station noise as diffusion v2 is (ratio of spread to error 0.96 for the quietest and 1.14--1.15 for the noisiest quarter of stations, 90\% coverage 0.90--0.92 in each); without it, its ensembles are slightly wider and heavier-tailed than diffusion v2's (spread to error 1.14--1.16 against 1.12--1.14, Tail 0.18 against 0.16), as if it allowed for a transient that might be present. Against the space--time Gaussian process it keeps the lead of diffusion v2: $-4.3$ and $-3.2\%$ in fast CRPS with the noise level, $-3.8$ and $-2.7\%$ without, better at 138--143 of the 154 stations (Table~\ref{tab:significance}). At 25 and 100~km, where diffusion v2 was not evaluated, diffusion v3 without the noise level scores as the first diffusion model (fast CRPS 0.280 and 0.291, against 0.281 and 0.291) and 3.7 and 2.8\% below the Gaussian process, better at 142 of the 154 stations at both densities. Regress v2 is as accurate as regress with the full network and 0.9\% worse in CRPS at 60~km.

\paragraph{Planted events are passed on} On the full network, diffusion v3 recovers a median of 52\% of the planted signal per target and event, against 36\% for diffusion v2 and 56\% for the space--time Gaussian process (Figure~\ref{fig:brho}). Where the footprint is two to four times the distance to the nearest station, it recovers 29\%, where diffusion v2 recovers nothing and the Gaussian process 14\%; where it is more than four times that distance, 82\%, against 72\% and 91\%. Over all networks it matches the Gaussian process where the footprint is more than four times the gap (56\% against 57\%) and exceeds it where it is two to four times the gap (13\% against 5\%). On the networks thinned to 25~km its median is 0.21--0.28 against 0.07--0.08 for diffusion v2 and 0.15--0.18 for the Gaussian process. For afterslip on the full network the median is 81\%, against 60\% for diffusion v2 and 90\% for the Gaussian process, and the shape of what is passed on is closer to the event (median correlation 0.92 against 0.89; 0.98 for the Gaussian process). Pooled by signal energy the gain is small (0.21 against 0.19, Table~\ref{tab:brho}), because the pooled fraction is dominated by a few targets within 25~km of an event: five of the 76 afterslip target--event pairs carry 70\% of the detrended signal energy, and at these every method recovers less than half. Shallow slow slip, smaller than the gaps between stations, remains invisible. The transfer is not bought with spurious signal: at held-out stations that an event leaves still, diffusion v3 adds a median of 2.0\% of the neighbours' peak displacement on the full network, the least of all methods (2.3\% for diffusion v2, 2.6\% for the Gaussian process), and during the events its CRPS of the planted truth and coverage are those of diffusion v2 (0.307 and 0.83, against 0.308 and 0.82). Regress v2 gains as much (median 0.46 against 0.32 for regress on the full network, 29\% where the footprint is two to four times the gap, and the largest energy-pooled fraction of afterslip, 0.37), so the gain comes from the training data rather than from the generative formulation.

\paragraph{Ridgecrest} The real offsets show the largest change (Table~\ref{tab:ridgecrest}). With the full network diffusion v3 recovers a median of 0.85 of the observed offsets, against 0.53 for diffusion v2, with the error of the Gaussian process (2.7~mm), and 81\% of the offset components lie inside its 90\% range, against 15\% for the Gaussian process. With the network thinned to 60~km it recovers 0.68--0.73 (with and without the noise level), where every method trained without transients recovers 0.09--0.24, and its median error is 3.5--3.6~mm, against 9.1--12.2~mm; its error is smaller than that of every method trained without transients at 11 of the 13 stations, and its generated offsets point in the observed directions. Regress v2 recovers 0.88--0.95 of the offsets but with larger errors (3.2 and 5.8~mm). Only TOWG, 15~km from the rupture with an offset of 0.55~m and its nearest network station 35~km away, is missed by every method. The contrast with the Gaussian process at 60~km explains the gain. Its regional covariance, fitted to ordinary days, decays within 5--21~km, so at 60~km spacing, with the nearest stations 14--47~km from the targets, it attributes most of a step shared by the neighbours to their own noise. The models trained with transients have learned that displacement can be coherent over tens to hundreds of kilometres, which the far field of a magnitude 7 earthquake is. The planted family contains such steps by design (up to 50~mm, footprints up to 300~km), but it was drawn without reference to Ridgecrest, whose days and near-field stations were excluded from training.

\clearpage
\section{Status and decisions}
\label{sec:status}

The following decisions departed from the initial plan, each on the evidence given above: minimum-separation thinning replaced random subsampling (Table~\ref{tab:thinning}); a fifth of the training samples retain the full network with no minimum distance, and the minimum distance on thinned networks is drawn from 1--100~km instead of 10--100~km, raising the share of training samples with a neighbour within 5~km from 1\% to 3\%; each window's mean and trend are removed from every series (Table~\ref{tab:drift}); and the classifier test was moved from the baseline stage to the model stage. The baseline results set two targets for the diffusion model, a short-period CRPS below 0.29 at 60~km density and inter-station correlations that decrease with distance as those of real stations do; the first trained model meets both. Its ensembles appeared too wide; three remedies were considered. Screening the samples as the observations were screened changes little and is retained as part of scoring generative methods. Classifier-free guidance narrows the ensembles only by exaggerating the shared signal and is not used. Removing the earthquake offsets from the training residuals is set aside, since the model is as wide on its own training distribution, earthquake steps or not, as on the validation period. The excess is instead explained by the station's unknown noise level, over which the model, like the baseline, hedges; pooled over stations its short-period intervals are calibrated, and a residual long-period excess of 10--15\% remains. Given the station's noise level as an optional input, the second diffusion model is calibrated across stations and loses most of the excess variability of the first. Against Gaussian generators built on a covariance fitted to the training years, the diffusion model keeps a lead of about 3\% in CRPS, but its reproduction of the decrease of correlation with distance is shared with kriging + noise and the space--time Gaussian process, and the comparison with IDW + noise overstated it. Experiment B showed that, without training on localised transients, the diffusion model transfers planted events no better than the near-uniform averages and less well than the space--time Gaussian process. Training with synthetic transients planted in the inputs, as planned (Section~\ref{sec:model}), raised the transfer of planted events to near that of the Gaussian process on the full network and above it where events are two to four times larger than the gaps between stations, at no cost in CRPS on ordinary days, and made the diffusion model the most accurate method for the real Ridgecrest offsets on a network as sparse as T\"urkiye's; it is the model carried forward. The largest afterslip signals, at stations close to the event, are still passed on only in part by every method. The remaining work is the classifier test, the test period, and the T\"urkiye data cube and experiment C.

\section*{Acknowledgements}

We thank Dr.\ Bülent Bolat (Department of Electronics and Communications Engineering, Yıldız Technical University) for his guidance and support, and İsmail Haydar İmamoğlu for his support.

\bibliographystyle{elsarticle-harv}
\bibliography{references}

\end{document}
